\chapter*{附录\quad 易错表述与更正}
\addcontentsline{toc}{chapter}{附录\texorpdfstring{\quad}{ }易错表述与更正}
\markboth{附录\quad 易错表述与更正}{}
下表汇总学习中容易出现的不准确写法、记号或数据转录错误，并给出本书采用的写法，便于复习时对照。“位置”为本书的节号。
{\small\renewcommand{\arraystretch}{1.25}\setlength{\tabcolsep}{4pt}
\begin{longtable}{>{\raggedright\arraybackslash}p{1.25cm}>{\raggedright\arraybackslash}p{4.8cm}>{\raggedright\arraybackslash}p{4.55cm}>{\raggedright\arraybackslash}p{3.0cm}}
\toprule 位置&易错写法或表述&本书采用&说明\\\midrule\endfirsthead
\toprule 位置&易错写法或表述&本书采用&说明\\\midrule\endhead
\bottomrule\endfoot
1.1.1&组成表随机部分写作 $\mu=\widehat y$&$\mu_i=\E(Y_i\mid\bm x_i)$；拟合值 $\widehat\mu_i=g^{-1}(\bm x_i^\top\widehat{\bm\beta})$&区分总体量与估计值\\
1.1.1&“systemic component”&systematic component&拼写\\
1.1.2&Cox回归与GLM并列&注明是半参数生存回归模型&不属于标准GLM\\
1.2&表7取值范围 $(0,n)$、$(0,\infty)$&$\{0,1,\ldots,n\}$、$\{0,1,2,\ldots\}$；区分二项计数与比例&离散分布含端点\\
1.3.3&对数似然首项 $-\frac12\log(2\pi\sigma^2)$&$-\frac n2\log(2\pi\sigma^2)$&$n$ 个观察之和\\
1.3.3&“$f(\bm\beta,\sigma^2\mid\bm y)$：给定样本，参数的一个函数”&删去，只保留 $L(\bm\beta,\sigma^2\mid\bm y)=f(\bm y\mid\bm\beta,\sigma^2)$&似然不是参数的密度\\
1.3.3&$\bm X$ 为“$n(m+1)$ 的矩阵”&$n\times(m+1)$&\\
1.4&公式写作 \texttt{Y=X1+X2+X3}；\texttt{family=gaussion}&\texttt{Y \textasciitilde{} X1+X2+X3}；\texttt{gaussian}&公式符号与拼写\\
1.1，2.1.1&图横轴记为 $y$（logit值）&$\eta$&线性预测量不是响应\\
2.1.4&$H_0:\widehat\beta_i=0$&$H_0:\beta_i=0$&假设针对总体参数\\
2.1.4&比分检验批注：$U(\theta)$ 是“均数、对数似然函数”&$U$ 为得分函数，$I$ 为Fisher信息&\\
2.1.4&“Wald检验未考虑各因素间的综合作用”&删去&多变量Wald检验同样是调整后的\\
2.1.4&“似然比检验结果最可靠”&三者渐近等价，小样本时似然比检验通常更可靠&\\
2.1.5&“全模型（full model）”&饱和模型（saturated model）&\\
2.1.5&$\lambda$“越小，拟合优度越好”&$\lambda\ge1$，越接近1拟合越好&\\
2.1.5&广义 $\chi^2$ 分母 $\Var(\widehat y)$&模型方差 $n_g\widehat\pi_g(1-\widehat\pi_g)$&\\
2.1.5&Pseudo $R^2$ 未注明定义&McFadden定义&\\
2.2.2&条件概率中的比值记作 $\OR_{B_i}/\OR_{A_i}$；未写对子截距&odds之比；先引入 $\alpha_i$ 再消去&\\
2.2.2&表5批注：“用药天数越长，不良反应概率越大”&天数每增加1天，优势乘0.478&与系数符号相反\\
2.2.2&“合并药物每增加1个，发生概率增加0.012”&优势乘1.444&0.012是 $P$ 值\\
2.2.3&批注“暴露者 $X=0$”&非暴露者 $X=0$&\\
3.1.2&表7（分化2、染色$-$、泡状细胞癌）36&26&用26才能精确复现表8\\
3.1.2&表8末列“RR”&RRR（相对风险比）&\\
3.1.3&批注“分别比较：$P_1+P_3=1$，$P_2+P_3=1$”&分别拟合的是条件概率 $P_1/(P_1+P_3)$&\\
3.2&“等比例风险假设/假定”&比例优势假设&与生存分析的比例风险区分\\
3.2&“从1级升到2级、2级升到3级作用相同”&任一分界点两侧的OR相同&前者是相邻类别模型\\
3.2&不平行时可改用“有序Probit模型”&有序Probit同样要求平行&\\
3.2&平行线检验 $P>0.05$“认为满足假定”&未发现违背假设的充分证据&\\
3.2&例4：“提高一个及一个以上等级的可能性增加0.89倍”&任一分界点上“等级更高”的优势乘以1.89&\\
3.3.1&多项Logit“4个级别是相互独立的类别”&不使用类别顺序&\\
3.3.1&累积Logit“能回答病变进展速度”&横断面资料只能描述严重程度的关联&\\
3.3.1&“本案例不适合用多项Logit模型”&多项Logit仍有效；比例优势成立时累积Logit更有效率&\\
3.3.2&2848个研究对象；1439个“因痴呆死亡”的病例&1439对共2878人；结局为新发痴呆&$1224+140+60+15=1439$\\
3.3.2&研究设计“病例对照研究”&双生子队列中的共同双生子配对分析&\\
3.3.2&TABLE 1：Male 1804；$<8$ years 131,46&18,074；13,146&按百分比与合计核对\\
3.3.2&OR 2.36（95\%CI 1.74--2.30）&1.74--3.20&与TABLE 3一致\\
3.4&样本量“变量个数的20倍”&事件数为参数个数的10--20倍&\\
3.6&\texttt{require(mass)}&\texttt{library(MASS)}；响应为有序因子&包名区分大小写\\
3.6&SAS有序模型 \texttt{link=glogit}&\texttt{link=clogit}；\texttt{glogit} 用于无序多分类&\\
3.6&PROC LOGISTIC未指定事件水平&\texttt{y(event='1')} 或 \texttt{descending}&否则系数符号相反\\
4.1.6&“假定观察对象在各个时段的生存事件独立”&由条件概率的乘法公式得到连乘&不需要独立性\\
4.1.6&$S(t)=$ 存活例数/观察总例数&仅为无删失时的估计&\\
4.1.7&有的说法称 $F$ 为“累积的风险函数”&$F$ 为累积发生概率；$H=-\ln S$ 为累积风险&\\
4.2.1&“对数变换后 $G(t_i)$ 就服从正态分布”&近似服从正态；反变换递减，上下限互换&\\
4.2.2&KM法“一般用于观察对象数目较少的资料”&删去&适用于有确切时间的资料\\
4.2.2&例3中167未标“$+$”&$167^+$&与图9、表2、表5一致\\
4.2.2&表4（续）30月：$0.5401\times0.7500=0.3240$；62月：$0.2160\times0$&$0.4321\times0.7500=0.3241$；$0.1080\times0$&\\
4.3.1&表5第3行为删失时点 $t=3$（$T_{1i}=0.6071$）&只列事件时点&$d_i=0$ 时理论数为0\\
4.3.1&表5合计死亡数22&21&$11+10$\\
4.3.1&精确法：两组的 $(A_g-T_g)^2$ 各除以6.7578后相加&$(A_1-T_1)^2/\sum V_{1i}$，$V_1=3.378$&数值6.772不变\\
4.3.1&“分母经方差校正，卡方变大，更容易拒绝 $H_0$”&近似公式本身偏保守&\\
4.3.2&“Breslow针对短期效应，log-rank针对长期效应”&log-rank等权，Breslow以风险人数加权&\\
4.3.4&“各组生存曲线不能交叉，交叉提示混杂”&交叉使效能下降、提示非比例风险&\\
5.1.2&“风险比（risk ratio, RR；或hazard ratio, HR）”；$\RR_j=\exp(\beta_j)$&HR；与固定时点的RR区分&\\
5.1.2&风险集：“在 $t_k$ 时刻及之前仍未发生事件且未删失”&$t_k$ 前一瞬间仍在观察（$\widetilde T\ge t_k$）&\\
5.1.2&表8第5行：id 8，PI 4.739，$S(t)$ 0.767&id 5，PI 6.174，$S(t)$ 0.899&与表2及同PI各行一致\\
5.1.2&“肿瘤分级每增加1级，死亡险变为5.357倍”&4.816倍&$\exp(1.572)$\\
5.2.1&累积风险曲线“应大致平行或等距”&应成比例；$\ln H$（log--log图）平行&\\
5.2.2&外在时依协变量：“取值不随时间改变，HR随时间改变”&那是时变系数 $\beta(t)$；核辐射例归入时变系数&\\
5.2.2&“在Cox回归中，RR和HR是一样的”&删去&\\
5.2.3&例8：服药的标准差0.8，$n\approx135$&0/1变量标准差$\le0.5$；$1:1$ 分配时 $D=241$，$n=344$&\\
5.2.4&图2“复发累积风险率”&累积发生率（$1-\mathrm{KM}$ 与CIF）&\\
5.3.4&批注“高强度运动反而保护效果减弱”&按图读出HR；VPA约150 min/wk后趋平&\\
5.5&\texttt{surv(time,status)}&\texttt{Surv(time, status)}&\\
5.5&Stata代码以“\#”作注释&“//”&\\
6.1&按收入合并后白人“少”为160；表头“家庭收入”&106；表头“肤色”&$67+36+3$\\
6.1&癌症数据合并表：9050、14995&9005、14955&$9000+5$\\
6.2.1&直接令 $u_1(i)=\ln p_{i\cdot}$ 等；“四个独立参数，四个格子”&按均值分解定义；固定总数时3个自由概率参数&零和约束\\
6.2.2&$u_{23}$：“变量2第 $j$ 级与变量2第 $k$ 级”&变量3第 $k$ 级&\\
6.2.2&批注“性别与自杀方式有关”&删去&与本例无关\\
6.2.2&$i$：人种；$j$：分娩次数；$k$：收入差异&$j$：收入；$k$：分娩次数&与表2、表3一致\\
6.2.3&参数表：有色$\times$低收入、有色$\times$高收入的 $P$ 值互换；两行 $u_{23}$ 标签互换&已更正，并注明参照水平&\\
6.2.4&“$P>0.05$ 说明模型符合实际，模型好”&证据不足以否定模型&\\
6.2.5&$R_{\mathrm{std}}=(O-E)/\sqrt E\sim N(0,1)$&Pearson残差方差小于1；调整残差近似 $N(0,1)$&\\
6.2.5&批注“最大4.6超过2，两两关联模型不太好”&该残差表来自完全独立模型&\\
6.2.6&“当且仅当第三个变量独立于其中一个或两个时”&充分条件&\\
6.2.6&“交互效应均不为0，合并将歪曲变量间关系”&一般会改变，但不是必然&\\
6.3&“5254例”&5251例&表中合计\\
6.3&模型拟合表（df 40、38、36、34、32）&按资料复算（最大df 18）&\\
6.3&Logistic回归及交互项结果（如X1:Y1 OR 1.25，$P=0.01$）&按资料复算，未见交互作用&\\
6.3&“Logistic无法分析自变量与结局的交互”“通常需要个体资料”&对应关系见补充推导；分组资料同样适用&\\
6.5&SAS、Stata代码以“\#”作注释&“/* */”、“//”&\\
7.1&“$\log(\widehat Y)$ 可由线性组合拟合”&$\log\E(Y\mid\bm x)=\bm x^\top\bm\beta$&\\
7.1.2&$\ln P_{ij}=\ln(\mu_{ij}/n_{ij})$&$\ln r_{ij}$&发生率不是概率\\
7.1.2&表5：$G^2$ 167.607、7.862、133.976、2.493&167.034、7.810、133.450、2.475&R复算，结论不变\\
7.1.2&表6：系数0.35、1.17、1.90、3.40、4.10及相应RR、区间&0.409、1.311、2.140、3.020、3.790&与同表SE、Wald一致\\
7.1.2&表7单位“‰”&每千人年&\\
7.2.1&负二项均数 $(1+\kappa)\lambda$，方差 $(1+\kappa)\lambda(1+\lambda)$&均数 $\mu$，方差 $\mu+\alpha\mu^2$（NB2）&\\
7.2.2&过离散 $\chi^2$：模型C 27.1824、模型D 32.8259（df $=1$）&58人分析集为30.50、12.74；边界检验 $P$ 值减半&原值来自59人\\
7.2.2&“负二项不容易犯假阳性错误，结果更保守”&Poisson低估方差；NB修正了方差&\\
7.2.4&主分析“至少14天”、敏感性分析“至少7天”；人年 $\times1000$&主分析至少7天、敏感性分析至少14天；人日 $\times100000$&与原文方法一致\\
7.2.5&删除就诊次数为0者以满足等离散&撤销；全样本用负二项，或零截断模型&\\
7.3&\texttt{offset=} 留空；用qcc包检验 $Y$ 本身的过离散&\texttt{offset(log(t))}；以模型残差检查过离散&\\
7.3&SAS \texttt{model y/n=... offset=ln}&\texttt{model y=... offset=ln\_t}&Poisson响应为计数\\
7.3/7.4&Stata \texttt{exposure(offset)}、\texttt{offset(offset)}&\texttt{exposure(t)} 或 \texttt{offset(ln\_t)}&\\
8&各章对应表述&按以上各条同步更正&\\
\end{longtable}}

\section*{资料局限说明}
\begin{itemize}
\item 第5章例5：原资料只给出30例中的11例记录，表9无法复算；表8的PI列所用系数与表9略有出入，8号患者记录与5号完全相同，原因不明。relapse没有记录复发时间，无法检查是否引入了未来信息。
\item 第7章癫痫例：本书所用资料与公开的Thall--Vail资料（R \texttt{MASS::epil}）有7行在发作次数或年龄上不同；本书按所列资料计算，可复现原分析的模型C、D。
\item 第7章职业暴露例：$G^2$ 与原资料相差0.02--0.57，原因不明；Pearson $\chi^2$、系数标准误、Wald $\chi^2$、预测值和残差均与原资料一致。
\item 第3章案例1、第7章儿童发热例的结果表为示意数据，没有原始资料，未作复算。
\end{itemize}
